\chapter[Introdução]{Introdução}
\label{cap_1}

\section{Caracterização do Problema}

A segmentação de imagens é uma tarefa fundamental em Visão Computacional, com aplicações em análise de imagens médicas, diagnóstico assistido, monitoramento e inspeção automatizada. Os avanços do Aprendizado Profundo (\ac{DL}) ampliaram o desempenho dos modelos, mas o obstáculo que motiva esta dissertação permaneceu: redes de segmentação precisam de grandes volumes de imagens anotadas pixel a pixel, e produzir essa anotação exige tempo de um especialista. Em domínios biomédicos o problema é agudo, porque o especialista é escasso e a anotação precisa acompanhar fronteiras que nem sempre são evidentes.

Uma linha de trabalho recente propõe contornar esse custo por outro caminho. Em vez de ajustar milhões de parâmetros por retropropagação a partir de muitos exemplos, o \ac{FLIM} \cite{souza2020flim,benato2020cnnmarkers} constrói os filtros da rede \emph{diretamente} a partir de traços que o especialista desenha sobre poucas imagens. Os trechos de imagem sob esses traços são agrupados, e os centróides resultantes passam a ser os núcleos da convolução. Não há função de perda, não há otimizador e não há gradiente. O treinamento custa segundos e opera com três a cinco imagens anotadas.

O trabalho de referência reproduzido nesta dissertação \cite{soares2026flim} leva essa ideia a um conjunto de decodificadores adaptativos e reporta resultados competitivos com redes treinadas por retropropagação. Mas há nele um ponto que merece atenção. A escolha das poucas imagens que o especialista vai anotar é feita por um procedimento iterativo que, a cada passo, mede o desempenho do modelo em todas as imagens disponíveis e escolhe aquela em que ele vai pior. Medir o desempenho em uma imagem exige a anotação de referência dela. Escolher três imagens entre as 848 do conjunto exige, portanto, que as 848 já estejam anotadas.

Os próprios autores declaram que a abordagem de seleção é supervisionada. A consequência, contudo, é mais severa do que a palavra sugere: um procedimento que precisa da anotação de todo o conjunto para decidir o que anotar devolve exatamente o custo que o \ac{FLIM} foi criado para evitar.

Daí a pergunta que organiza este trabalho:

\begin{quote}
Quanto do benefício daquele procedimento se recupera com critérios de seleção que \textbf{não consultam anotação alguma}?
\end{quote}

A pergunta pertence ao campo do Aprendizado Ativo (\ac{AL}) \cite{settles2009active}, que estuda precisamente como escolher o que anotar usando apenas o modelo corrente. A literatura de \ac{AL} é extensa e madura, e a expectativa razoável seria que seus critérios transferissem para o \ac{FLIM} com ajustes menores. O que esta dissertação mostra é que não transferem, e, mais importante do que constatar isso, \emph{por que} não transferem.

\section{Motivações e Desafios}

O interesse da pergunta não está apenas em economizar anotação. Está em que o \ac{FLIM} é um modelo com propriedades incomuns, e essas propriedades mudam o que cada técnica significa.

\textbf{Não há gradiente.} A maior parte do arsenal de Aprendizado Contínuo (\ac{CL}) consolidado na literatura, da consolidação elástica de pesos \cite{kirkpatrick2017overcoming} à destilação de conhecimento \cite{hinton2015distilling} e ao ensaio por exemplares \cite{rebuffi2017icarl}, atua sobre o deslocamento de pesos causado pelo otimizador. Em um modelo sem otimizador, esses métodos não têm sobre o que agir. Isso não significa que não haja esquecimento: significa que, se houver, ele tem outra forma, e descrevê-la é parte do problema.

\textbf{A anotação não é monotonicamente boa.} Em uma rede treinada por retropropagação, acrescentar um exemplo rotulado correto raramente piora o modelo. No \ac{FLIM}, acrescentar um marcador amplia o conjunto de trechos candidatos a filtro, e o agrupamento que os reduz ao número de canais disponíveis é refeito. Os filtros em uso se deslocam, e como o banco é compartilhado por todas as imagens, a perturbação é global. Anotar mais pode piorar.

\textbf{O que se anota é um traço, não uma máscara.} A unidade de anotação do \ac{FLIM} não é a imagem nem o pixel isolado: é um traço, com posição, extensão e proporção entre objeto e fundo. Isso abre uma pergunta que a formulação usual do \ac{AL} não faz. A literatura pergunta \emph{quais imagens} anotar; aqui é possível perguntar também \emph{onde} anotar dentro de uma imagem já escolhida, mantendo fixo o orçamento de pixels marcados.

\textbf{O especialista está no laço.} A anotação por traços é rápida o bastante para ser feita em tempo real, o que torna natural acoplar Aprendizado Interativo (\ac{IL}) \cite{xu2016deep_arxiv,sofiiuk2020fbrs_arxiv} ao método, com o modelo propondo onde a dúvida está e o especialista respondendo.

Os três componentes, portanto, não são acessórios somados a um modelo pronto. Cada um deles interage com uma propriedade estrutural do \ac{FLIM}, e é essa interação que a dissertação investiga.

\section{Objetivos}

\subsection{Objetivo Geral}

Investigar se, e sob quais condições, técnicas de \ac{AL}, \ac{CL} e \ac{IL} transferem para o \ac{FLIM}, caracterizando com medição o mecanismo que explica o comportamento observado e propondo a intervenção que esse mecanismo indica.

A formulação é deliberada em dois aspectos. Primeiro, ela não promete redução de custo de anotação, porque o trabalho mede número de interações e qualidade de segmentação, e não tempo de especialista. Segundo, ela admite resposta negativa: estabelecer que uma família de técnicas não transfere, acompanhada do mecanismo que explica a falha, é resultado de pesquisa.

\subsection{Objetivos Específicos}

\begin{itemize}

\item Reproduzir o trabalho de referência \cite{soares2026flim} em condições controladas, estabelecendo a linha de base sobre a qual qualquer comparação se apoia.

\item Avaliar critérios de \ac{AL} de incerteza, de diversidade e híbridos na seleção de imagens, usando apenas informação disponível sem anotação, e comparando contra dois referenciais: o sorteio aleatório, como piso, e o procedimento supervisionado do trabalho de referência, como teto.

\item Estender ao nível de região os critérios de diversidade, que na literatura são de imagem por construção, e avaliar a seleção de \emph{onde} anotar com orçamento de pixels fixo.

\item Medir a margem disponível para a seleção de região, isto é, quanto separa a melhor posição de anotação examinada de uma posição sorteada, e quanto dessa margem os critérios efetivamente capturam.

\item Diagnosticar onde os critérios caem quando erram, e distinguir se a limitação está no escore ou no conjunto de candidatos que ele recebe.

\item Propor e avaliar uma intervenção sobre o gerador de candidatos, mantendo escore e orçamento fixos, de modo a isolar o efeito do gerador.

\item Caracterizar a forma que o esquecimento assume em um modelo sem gradiente, propor um mecanismo de retenção compatível com essa forma, e testá-lo.

\item Medir o esforço de interação em um laço de \ac{IL} com usuário simulado, pela métrica de número de cliques.

\item Desenvolver uma plataforma local que integre os três componentes em um fluxo único, permitindo observar o comportamento do sistema com um especialista real no laço.

\item Conduzir toda a avaliação sob protocolo declarado antes da execução, com teste pareado, unidade de análise explícita e correção para múltiplas comparações.

\end{itemize}

\section{Contribuições}

As contribuições abaixo são as efetivamente obtidas, com a medição que as sustenta indicada entre parênteses. O Capítulo~\ref{cap_4} apresenta cada uma com a tabela correspondente.

\textbf{Primeira, a caracterização do motivo pelo qual as técnicas não transferem, com cada elo do argumento medido.} Nenhum critério de seleção de imagem supera o sorteio aleatório de forma que sobreviva à correção para múltiplas comparações, e o procedimento supervisionado do trabalho de referência tampouco se destaca. Mas existe margem considerável a capturar na seleção de \emph{região}: anotar a melhor região examinada em vez de uma sorteada vale entre 0,1351 e 0,3089 de $F_\beta$, nas seis combinações de conjunto de dados e decodificador, com significância em todas. Os critérios capturam apenas parte dela, e a parte não capturada, entre 0,0839 e 0,2667, é o custo de o critério escolher o lugar errado.

\textbf{Segunda, a identificação do gargalo e a intervenção que ele indica.} A limitação não está no escore. Regiões que atravessam a fronteira do objeto representam de 2\% a 12\% dos candidatos produzidos por superpixels, e um critério apenas ordena a lista que recebe. Substituir esse gerador por discos centrados no contorno predito, mantendo escore e orçamento fixos, elevou o ganho de uma anotação adicional em 0,0772 de $F_\beta$ no conjunto de parasitas, com vantagem nas dez repetições e sobrevivendo à correção. O comportamento nos demais conjuntos é compatível com o mecanismo, incluindo um resultado nulo previsto antes da execução.

\textbf{Terceira, a descrição do esquecimento em um modelo sem gradiente.} O fenômeno existe no \ac{FLIM} e assume forma própria: a anotação adicional amplia o conjunto de candidatos a filtro, o agrupamento é refeito e os filtros em uso se deslocam, com perturbação global porque o banco é compartilhado. Os três grupos de métodos consolidados atuam sobre deslocamento de pesos por gradiente e não se aplicam. O mecanismo de retenção proposto resolve o problema a que se destina, verificado por teste, mas não se converteu em ganho mensurável em nenhum regime avaliado.

\textbf{Quarta, de natureza metodológica, com alcance além deste trabalho.} O treinamento do codificador não era reprodutível entre execuções quando o número de candidatos ultrapassava o limite da arquitetura, e a causa foi localizada na interação entre o agrupamento e a ordem de redução em bibliotecas de álgebra linear com múltiplas linhas de execução. Verificou-se também que parâmetros de pós-processamento publicados não transferem entre domínios. E o protocolo de pré-registro adotado capturou quatro falsos positivos que teriam entrado no texto como achados.

\textbf{Quinta, a plataforma interativa.} Uma aplicação local integra os três componentes em um fluxo único, no qual o \ac{AL} propõe a região, o especialista a rotula, e a plasticidade do banco de filtros pode ser alterada entre interações. Ela permite observar em tempo real o fenômeno descrito aqui, incluindo o caso em que anotações sucessivas não melhoram o modelo.

\section{Questões e Hipóteses da Pesquisa}

As hipóteses abaixo foram declaradas antes da execução das campanhas correspondentes, com o critério de falseamento escrito junto. O veredito de cada uma consta ao lado, e três delas foram refutadas pelos próprios critérios que as acompanhavam. Hipótese refutada por medição é resultado, e não falha de planejamento; registrá-las é o que torna as demais confiáveis.

\begin{itemize}

\item \textbf{Q1.} Quanto do benefício do procedimento supervisionado de seleção se recupera com critérios de \ac{AL} que não consultam anotação?
\begin{itemize}
\item \textbf{H1.} Ao menos um dos critérios avaliados supera o sorteio aleatório na seleção de imagens. \textbf{Refutada}: em 96 testes, nenhum sobrevive à correção para múltiplas comparações, e o próprio procedimento supervisionado também não se destaca.
\end{itemize}

\item \textbf{Q2.} Se a escolha de quais imagens anotar não é a alavanca, a escolha de \emph{onde} anotar, com orçamento de pixels fixo, é?
\begin{itemize}
\item \textbf{H2.} Existe margem mensurável entre a melhor posição de anotação e uma posição sorteada. \textbf{Sustentada}: entre 0,1351 e 0,3089 de $F_\beta$, em 6 de 6 células, com significância em todas.
\item \textbf{H2b.} Os critérios de \ac{AL} capturam essa margem. \textbf{Parcialmente refutada}: capturam parte dela, e a lacuna remanescente, de 0,084 a 0,267, é significativa em 4 de 6 células.
\end{itemize}

\item \textbf{Q3.} A limitação está no escore que ordena os candidatos ou no procedimento que gera a lista de candidatos?
\begin{itemize}
\item \textbf{H3.} Está na lista. \textbf{Sustentada}: substituir o gerador, mantendo escore e orçamento fixos, produz ganho de 0,0772 de $F_\beta$ no conjunto de parasitas, em 10 de 10 repetições, sobrevivendo à correção; e o resultado nulo nos demais conjuntos foi previsto pelo mecanismo antes da execução.
\end{itemize}

\item \textbf{Q4.} O esquecimento catastrófico ocorre no \ac{FLIM}, e os métodos consolidados de \ac{CL} se aplicam a ele?
\begin{itemize}
\item \textbf{H4.} Ocorre, em forma distinta, e os métodos consolidados não se aplicam. \textbf{Sustentada} quanto à descrição do fenômeno.
\item \textbf{H4b.} Reter o banco de filtros preserva qualidade ao longo das interações. \textbf{Refutada} pelo critério declarado antes da execução.
\end{itemize}

\item \textbf{Q5.} Congelar o banco de filtros reduz o número de interações necessárias no laço de \ac{IL}?
\begin{itemize}
\item \textbf{H5.} Reduz. \textbf{Refutada}: o efeito não é detectável em nenhum dos cortes limpos, e um achado aparente na análise agregada desapareceu quando os conjuntos de dados foram separados.
\end{itemize}

\end{itemize}

\section{Organização do texto}

O Capítulo~\ref{cap_2} apresenta os conceitos necessários e os trabalhos relacionados, cobrindo segmentação, superpixels, Aprendizado Ativo, Contínuo e Interativo, e o próprio \ac{FLIM}, com atenção a por que a combinação entre eles não é imediata. O Capítulo~\ref{cap_3} descreve o método em detalhe, incluindo a formulação do \ac{FLIM}, as métricas, cada critério de \ac{AL} avaliado com a justificativa de sua inclusão, o mecanismo de retenção proposto, o usuário simulado do laço interativo, a plataforma, e o protocolo experimental. O Capítulo~\ref{cap_4} apresenta os resultados na ordem do raciocínio da investigação, e não na ordem cronológica dos experimentos. O Capítulo~\ref{cap_5} responde às questões, enuncia as contribuições e as limitações, e aponta os desdobramentos.

\section{Publicações}

Encontra-se em elaboração um artigo científico para divulgação dos resultados relativos à intervenção no gerador de candidatos e ao diagnóstico que a motiva, com submissão prevista para conferência da área. Os resultados de natureza metodológica, em particular o protocolo de pré-registro adotado e os defeitos de reprodutibilidade identificados e corrigidos, são candidatos a submissão separada, por terem alcance além do método estudado aqui.
